% Almacén spine_dw.db (sql/dw.sql): constelación de tres hechos con dimensiones compartidas;
% jerarquía anatómica normalizada (copo de nieve). Compilar con XeLaTeX.
\documentclass[border=6pt]{standalone}
\usepackage{fontspec}
\usepackage[sfdefault]{FiraSans}
\usepackage{tikz}
\usetikzlibrary{positioning, arrows.meta, shapes.multipart, calc, backgrounds}
% Estilo compartido por los diagramas ER y del almacén (colores ITBA, iguales a defensa.tex).
\definecolor{itba}{HTML}{02568A}
\definecolor{itbaclaro}{HTML}{7FA9C8}
\definecolor{itbafondo}{HTML}{EAF1F7}
\definecolor{tinta}{HTML}{1F2933}
\definecolor{gris}{HTML}{8A949E}
\definecolor{alerta}{HTML}{C0392B}
\hyphenpenalty=10000 \exhyphenpenalty=10000   % sin división silábica dentro de las tablas
\tikzset{
  ent/.style={rectangle split, rectangle split parts=2, draw=itba, line width=0.8pt, rounded corners=2pt,
              rectangle split part fill={itba,white}, font=\scriptsize, align=left, inner sep=3pt,
              text width=#1, text=tinta},
  ent/.default=3.4cm,
  hecho/.style={ent=#1, rectangle split part fill={alerta,white}, draw=alerta},
  dimen/.style={ent=#1},
  linea/.style={draw=itba, line width=0.8pt},
  fina/.style={draw=itbaclaro, line width=0.7pt},
  card/.style={font=\footnotesize\bfseries, text=itba, inner sep=1.5pt},
  verbo/.style={font=\scriptsize\itshape, text=tinta, inner sep=1pt},
  nm/.style={draw=alerta, line width=0.8pt, dashed},
  nmtexto/.style={font=\scriptsize\bfseries, text=alerta},
}
\newcommand{\pk}[1]{\textbf{PK}~#1}
\newcommand{\fk}[1]{\textbf{FK}~#1}
\newcommand{\pkfk}[1]{\textbf{PK·FK}~#1}
\newcommand{\nulo}[1]{\textcolor{gris}{\textit{#1}}}
\newcommand{\fkn}[1]{\textbf{FK}~\nulo{#1}}
\newcommand{\llave}[1]{{\upshape\tiny\color{gris}#1}}
% \entidad[ancho]{nombre}{(x,y)}{atributos}
\newcommand{\entidad}[4][3.4cm]{\node[ent=#1] (#2) at #3 {\textcolor{white}{\textbf{\small #2}}\nodepart{two}#4};}
\newcommand{\dimension}[4][3.2cm]{\node[dimen=#1] (#2) at #3 {\textcolor{white}{\textbf{\small #2}}\nodepart{two}#4};}
\newcommand{\hechos}[4][3.6cm]{\node[hecho=#1] (#2) at #3 {\textcolor{white}{\textbf{\small #2}}\nodepart{two}#4};}
% \relacion{inicio}{fin}{verbo}{columna FK}: línea 1 (lado del padre, inicio) —— N (lado del hijo, fin)
\newcommand{\relacion}[4]{%
  \draw[linea] (#1) -- (#2)
    node[pos=0.07, above, card] {1} node[pos=0.93, above, card] {N}
    node[midway, above=1pt, verbo] {#3} node[midway, below=1pt, verbo, align=center] {\llave{#4}};}

\newcommand{\suma}[1]{\textbf{\color{alerta}Σ}~#1}
\newcommand{\prom}[1]{\textbf{\color{gris}x̄}~#1}
% \tabla{estilo=ancho}{nombre del nodo}{(x,y)}{título}{cuerpo}
\newcommand{\tabla}[5]{\node[#1] (#2) at #3 {\textcolor{white}{\textbf{\small #4}}\nodepart{two}#5};}

\begin{document}
\begin{tikzpicture}
  % ------------------------------------------------------------------ hechos (centro)
  \tabla{hecho=5.5cm}{fimg}{(8.1,4.35)}{fact\_imagen}{%
    \textit{Grano: una radiografía} · P1, P3, P4\\[1pt]
    \pkfk{imagen\_key}\\
    \fk{dataset\_key, diagnostico\_key}\\
    \fk{sexo\_key, grupo\_edad\_key}\\
    \fk{vista\_key, orientacion\_key}\\[1pt]
    \suma{imagenes, anotada, completa}\\
    \suma{vertebras\_anotadas, edad}}
  \tabla{hecho=5.5cm}{fvert}{(8.1,0.0)}{fact\_vertebra\_anotada}{%
    \textit{Grano: una vértebra en una imagen} · P1, P2, P3\\[1pt]
    \pkfk{imagen\_key, vertebra\_key}\\
    \fk{dataset\_key, diagnostico\_key}\\
    \fk{vista\_key, orientacion\_key}\\[1pt]
    \suma{anotaciones, es\_atipico, orden\_coherente}\\
    \prom{area\_norm, alto\_norm, aspecto, relleno,}\\
    \hspace*{0.9em}cy\_rel, inclinacion, excentricidad, …}
  \tabla{hecho=5.5cm}{fseg}{(8.1,-4.05)}{fact\_segmento}{%
    \textit{Grano: un segmento en una imagen} · P5\\[1pt]
    \pkfk{imagen\_key, segmento\_key}\\
    \fk{dataset\_key, diagnostico\_key, vista\_key}\\[1pt]
    \suma{segmentos}\\
    \prom{delta\_inclinacion, desplazamiento\_rel,}\\
    \hspace*{0.9em}separacion\_rel}

  % ------------------------------------------------------------------ dimensiones de un solo hecho (izquierda)
  \tabla{dimen=3.1cm}{dsexo}{(1.7,5.35)}{dim\_sexo}{\pk{sexo\_key}\\sexo {\color{gris}(F, M, sin dato)}}
  \tabla{dimen=3.1cm}{dedad}{(1.7,3.35)}{dim\_grupo\_edad}{\pk{grupo\_edad\_key}\\grupo, edad\_min, edad\_max}
  \tabla{dimen=3.1cm}{dvert}{(1.7,0.0)}{dim\_vertebra}{\pk{vertebra\_key}\\codigo, orden\_craneocaudal\\\fk{region\_key}}
  \tabla{dimen=3.1cm}{dseg}{(1.7,-4.05)}{dim\_segmento}{\pk{segmento\_key}\\nombre {\color{gris}(L1-L2 …)}\\\fk{vertebra\_sup\_key}\\\fk{vertebra\_inf\_key}}
  \tabla{dimen=1.9cm}{dreg}{(-2.15,0.0)}{dim\_region}{\pk{region\_key}\\nombre}

  % ------------------------------------------------------------------ dimensiones compartidas (derecha)
  \tabla{dimen=3.3cm}{dds}{(14.75,5.2)}{dim\_dataset}{\pk{dataset\_key}\\nombre, licencia, geometria}
  \tabla{dimen=3.3cm}{ddiag}{(14.75,3.0)}{dim\_diagnostico}{\pk{diagnostico\_key}\\nombre {\color{gris}(0 = sin dato)}}
  \tabla{dimen=3.3cm}{dvista}{(14.75,0.95)}{dim\_vista}{\pk{vista\_key}\\vista {\color{gris}(frente, perfil)}}
  \tabla{dimen=3.3cm}{dori}{(14.75,-1.1)}{dim\_orientacion}{\pk{orientacion\_key}\\orientacion {\color{gris}(0 = sin dato)}}
  \tabla{dimen=3.3cm}{dimg}{(14.75,-3.75)}{dim\_imagen}{\pk{imagen\_key}\\archivo, split, resolucion\\formato, canales, bits\\nivel\_listesis}

  % ------------------------------------------------------------------ dimensión (1) → hecho (N)
  \draw[linea] (dsexo.east) -- ([yshift=0.5cm]fimg.west);
  \draw[linea] (dedad.east) -- ([yshift=-0.5cm]fimg.west);
  \draw[linea] (dvert.east) -- (fvert.west);
  \draw[linea] (dseg.east) -- (fseg.west);
  \draw[linea] (dreg.east) -- (dvert.west);
  \draw[linea] (dvert.south) -- (dseg.north) node[midway, left, verbo, align=right] {superior e\\inferior};
  \foreach \d/\f in {dds/fimg, dds/fvert, dds/fseg, ddiag/fimg, ddiag/fvert, ddiag/fseg,
                     dvista/fimg, dvista/fvert, dvista/fseg, dori/fimg, dori/fvert,
                     dimg/fimg, dimg/fvert, dimg/fseg}
    \draw[fina] (\d.west) -- (\f.east);

  % ------------------------------------------------------------------ leyenda
  \node[anchor=north west, font=\scriptsize, text=tinta, align=left] at (-3.35,-6.3) {%
    \textcolor{alerta}{\textbf{rojo}}: tablas de hechos · \textcolor{itba}{\textbf{azul}}: dimensiones ·
    toda relación dimensión\,→\,hecho es \textbf{1\,:\,N} (cada fila de hechos referencia una fila de cada dimensión)\\
    \suma{} medida aditiva: se suma (\texttt{SUM}) para contar · \prom{} descriptor no aditivo: se promedia o se compara al consultar (nunca se guarda un promedio)\\
    izquierda: dimensiones de un solo hecho · derecha: dimensiones compartidas (conformadas) · región → vértebra → segmento: jerarquía normalizada (copo de nieve)};
\end{tikzpicture}
\end{document}
